\section*{الفصل \ref{ch:b2:series} --- المتتاليات والمتسلسلات}

\begin{solution}{exo:b2:series:1}
$\sum\frac{\cos n}{n}$: محك آبل مع $a_n = \frac1n$ ومع $b_n =
\cos n$، التي مجاميعها الجزئية محدودة (فهي الجزء الحقيقي لمجموع
هندسي، كما في \cref{ex:b2:series:sinn}): فهي متقاربة (لا \omterm{def:b2:series:def}{بإطلاق}،
بحيلة $\cos^2$ نفسها).

$\sum \frac{(-1)^n}{\ln n}$ (مع $n \geq 2$): محك المتناوبات،
$\frac{1}{\ln n} \downarrow 0$: فهي متقاربة؛ ولا \omterm{def:b2:series:def}{بإطلاق} ($\ln n
\leq n$).

$\sum \frac{(-1)^n}{n^{3/4} + \cos n}$: ننشر،
\[
\frac{(-1)^n}{n^{3/4} + \cos n}
= \frac{(-1)^n}{n^{3/4}}\cdot
\frac{1}{1 + \frac{\cos n}{n^{3/4}}}
= \frac{(-1)^n}{n^{3/4}}
- \frac{(-1)^n\cos n}{n^{3/2}} + O\Bigl(\frac{1}{n^{9/4}}\Bigr).
\]
المتسلسلة الأولى: متناوبة، فمتقاربة. والثانية: متقاربة
\omterm{def:b2:series:def}{بإطلاق} (بسلّم $\frac{1}{n^{3/2}}$). والثالثة: متقاربة
\omterm{def:b2:series:def}{بإطلاق}. والمحصلة: متقاربة.
\end{solution}

\begin{solution}{exo:b2:series:2}
\omterm{thm:b2:series:fubini}{جداء كوشي} للمتسلسلة $\sum_{m\geq0} z^m$ في نفسها (وكلتاهما متقاربة \omterm{def:b2:series:def}{بإطلاق}
من أجل $\abs z < 1$): فالمعامل القطري هو $c_k =
\sum_{m=0}^{k} 1 = k + 1$، ومنه
\[
\frac{1}{(1-z)^2} = \sum_{k\geq0} (k+1)z^k .
\]
وبالضرب في $z$ وإعادة الترقيم: $\sum_{n \geq 1} n z^n =
\frac{z}{(1-z)^2}$.
\end{solution}

\begin{solution}{exo:b2:series:3}
لتكن $B_n = \sum_{k \leq n} b_k \to B$. \omterm{thm:b2:series:abel}{جمع آبل} مع $a_k = k$:
\[
\sum_{k=1}^{n} k\,b_k = n B_n - \sum_{k=1}^{n-1} B_k
\quad\Longrightarrow\quad
\frac1n \sum_{k=1}^{n} k b_k = B_n - \frac{1}{n}\sum_{k=1}^{n-1}
B_k .
\]
وتؤول متوسطات تشيزارو للمتسلسلة المتقاربة $(B_k)$ إلى نهايتها $B$
(مجلد السنة الأولى)، ومن ثم يؤول الطرف الأيمن إلى $B - B = 0$.
\end{solution}

\begin{solution}{exo:b2:series:4}
إذا كان $\theta \in \pi\Z$: فكل الحدود تنعدم --- فهي متقاربة بداهةً.
ولنفترض $\theta \notin \pi\Z$.

$\alpha > 1$: متقاربة \omterm{def:b2:series:def}{بإطلاق} (بالهيمنة بالمقدار $n^{-\alpha}$).

$0 < \alpha \leq 1$: ينطبق محك آبل (إذ $a_n = n^{-\alpha}
\downarrow 0$؛ والمجاميع الجزئية للمقدار $\sin n\theta$ محدودة بالمقدار
$\frac{1}{\abs{\sin(\theta/2)}}$، وهو مجموع هندسي): فهي متقاربة. ولا
\omterm{def:b2:series:def}{بإطلاق}: إذ $\abs{\sin n\theta} \geq \sin^2 n\theta = \frac{1 -
\cos 2n\theta}{2}$، و$\sum \frac{1 - \cos 2n\theta}{2n^\alpha}$
متباعدة (لأن $\sum n^{-\alpha}$ متباعدة؛ و$\sum
\frac{\cos 2n\theta}{n^\alpha}$ متقاربة بمحك آبل حين $2\theta
\notin 2\pi\Z$؛ والحالة المستبعَدة $2\theta \in 2\pi\Z$ تعني
$\theta \in \pi\Z$، وقد عولجت). فهي نصف متقاربة.
\end{solution}

\begin{solution}{exo:b2:series:5}
لتكن $p_1, p_2, \dots$ الحدود غير السالبة في $(u_n)$ بترتيبها،
و$q_1, q_2, \dots$ الحدود السالبة. وتتباعد $\sum p_k$ و$\sum
q_k$ معًا: فلو تقاربت إحداهما، لساوت الأخرى
المتسلسلةَ المتقاربة $\sum u_n$ ناقصها، فتقاربت أيضًا --- وعندئذٍ
لتقاربت $\sum \abs{u_n} = \sum p_k - \sum q_k$، وهذا يناقض
الفرض. وأيضًا $u_n \to 0$ (لأن $\sum u_n$ تتقارب).

إعادة الترتيب الجشعة: خذ حدودًا موجبة $p_1, p_2, \dots$ إلى أن
يتجاوز المجموع الجاري $\ell$ أول مرة (وهذا ممكن لأن $\sum p_k =
+\infty$)؛ ثم حدودًا سالبة إلى أن ينزل المجموع تحت
$\ell$ أول مرة (وهذا ممكن لأن $\sum q_k = -\infty$)؛ وكرّر إلى الأبد (فكل
طور منته، ويُستعمل كل حدّ مرة واحدة بالضبط: فهي إعادة ترتيب
حقيقية). وبعد كل تبديل، تكون المسافة من المجموع الجاري
إلى $\ell$ لا تتجاوز آخر حدّ مستعمل؛ وبما أن الحدود المستعملة
عند التبديل رقم $m$ لها دليل $\to \infty$، وبما أن $u_n \to 0$،
فإن المجاميع الجارية تتقارب إلى $\ell$.
\end{solution}

\begin{solution}{exo:b2:series:6}
القابلية للجمع: المجاميع الجزئية المنتهية للمقدار
$\frac{\abs x^{m+n}}{m!n!}$ محدودة بالمقدار $\bigl(\sum_m
\frac{\abs x^m}{m!}\bigr)^2 = \eu^{2\abs x}$. والتجميع القطري
(\cref{thm:b2:series:fubini}):
\[
(\eu^{x})^2 = \sum_{m,n} \frac{x^{m+n}}{m!\,n!}
= \sum_{k=0}^{\infty} x^k \sum_{m+n=k} \frac{1}{m!\,n!}
= \sum_k \frac{x^k}{k!}\sum_{m=0}^{k}\binom km
= \sum_k \frac{(2x)^k}{k!} = \eu^{2x} .
\]
\end{solution}

\begin{solution}{exo:b2:series:7}
القابلية للجمع: محدودة بالمقدار $\zeta(2)^2$ بوصفها عائلة جداء
(بحجة \omterm{thm:b2:series:fubini}{جداء كوشي} في \cref{thm:b2:series:fubini}). والتطبيق
$(q, a, b) \mapsto (qa, qb)$، من الثلاثيات ذوات $\gcd(a, b) = 1$
إلى الأزواج $(m, n)$، تقابل (ضع $q = \gcd(m,n)$). وبتجميع
\omterm{def:b2:series:summable}{العائلة القابلة للجمع} تبعًا لذلك (وهو تقسيم لمجموعة الأدلة ---
مشروع من أجل العائلات \omterm{def:b2:series:summable}{القابلة للجمع} حسب
\cref{thm:b2:series:rearrangement}/\cref{thm:b2:series:fubini}
مطبَّقة على التقسيم إلى أصناف \omterm{def:b2:structures:countable}{قابلة للعد}):
\[
\zeta(2)^2 = \sum_{q} \sum_{\gcd(a,b)=1} \frac{1}{q^4 a^2 b^2}
= \zeta(4)\, S,
\qquad\text{ومنه}\qquad
S = \frac{\zeta(2)^2}{\zeta(4)} .
\]
(ومع القيمتين $\zeta(2) = \frac{\pi^2}{6}$ و
$\zeta(4) = \frac{\pi^4}{90}$ من
\cref{ch:b2:fourier}: $S = \frac{5}{2}$.)
\end{solution}

\begin{solution}{exo:b2:series:8}
لتكن $A = \sum a_n$ (\omterm{def:b2:series:def}{بإطلاق})، و$B_m = \sum_{k\leq m} b_k \to B$،
محدودة بالمقدار $M$. عندئذٍ
\[
C_N = \sum_{k=0}^{N} c_k = \sum_{n=0}^{N} a_n B_{N-n}
\]
(بالجمع حسب دليل $a$). ونكتب
\[
C_N - AB = \sum_{n=0}^{N} a_n (B_{N-n} - B) - B\sum_{n > N} a_n .
\]
ويؤول الحدّ الأخير إلى $0$. ونفصل المجموع عند $n = \lfloor N/2
\rfloor$: فمن أجل $n \leq N/2$، $N - n \geq N/2$، ومنه $\abs{B_{N-n} - B}
\leq \varepsilon_N := \sup_{m \geq N/2}\abs{B_m - B} \to 0$، ويكون
هذا الجزء $\leq \varepsilon_N \sum\abs{a_n}$؛ ومن أجل $n > N/2$،
$\abs{B_{N-n} - B} \leq 2M$، ويكون هذا الجزء $\leq 2M \sum_{n >
N/2} \abs{a_n} \to 0$. ومنه $C_N \to AB$.
\end{solution}

\begin{solution}{exo:b2:series:9}
نأخذ $u_n = 2^{-n} e_n$ (حيث $e_n$ المتتالية التي فيها $1$ واحد في
الموضع $n$). عندئذٍ $\sum \norm{u_n}_\infty = \sum 2^{-n} <
\infty$: فهي متقاربة \omterm{def:b2:series:def}{بإطلاق}. لكن المجاميع الجزئية $S_N =
(1, \tfrac12, \dots, 2^{-N}, 0, \dots)$ كان يجب أن تتقارب إلى
المتتالية $(2^{-n})_n$، وهي \emph{ليست} منعدمة ابتداءً من رتبة ما:
فهي خارج $E$. وداخل $E$، تكون $(S_N)$ كوشي بلا نهاية (إذ إن $\norm{S_N
- x}_\infty \geq 2^{-N-1}$ لا يفيد أي متتالية $x$ منعدمة ابتداءً من رتبة ما:
إذ من أجل أي $x \in E$ منعدمة بعد الرتبة $K$، يكون $\norm{S_N - x} \geq
2^{-K-1}$ من أجل $N > K$): فالمتسلسلة لا تتقارب في $E$.
والتمام هو بالضبط ما تحتاجه \cref{thm:b2:nvs:absoluteconvergence}.
\end{solution}

\begin{solution}{exo:b2:series:10}
كلٌّ من $\arctan(n+1) - \arctan n$ و
$\arctan\frac{1}{n^2+n+1}$ يقع في $\intoo{0}{\frac\pi2}$، وتعطي
صيغة جمع الظلال أن
\[
\tan\bigl(\arctan(n{+}1) - \arctan n\bigr)
= \frac{(n+1) - n}{1 + n(n+1)} = \frac{1}{n^2 + n + 1} :
\]
فالظلان متساويان في فترة يكون فيها $\tan$ متباينًا، ومن ثم تصحّ
المتطابقة. وبالتلسكوب،
\[
\sum_{n=1}^{N}\arctan\frac{1}{n^2+n+1}
= \arctan(N{+}1) - \arctan 1 \xrightarrow[N\to\infty]{}
\frac\pi2 - \frac\pi4 = \frac\pi4 .
\]

\end{solution}

\begin{solution}{exo:b2:series:11}
الأولى: $\sqrt{n+1} - \sqrt n = \frac{1}{\sqrt{n+1} + \sqrt n}
\sim \frac{1}{2\sqrt n}$، ومن ثم فالحدود $\sim
2^{-\alpha}n^{-\alpha/2}$: فالتقارب إذا وفقط إذا $\frac\alpha2 > 1$،
أي $\alpha > 2$. والثانية: $1 - \cos\frac1n \sim
\frac{1}{2n^2}$: فتتقارب. والثالثة: $\bigl(1 + \frac1n\bigr)^n =
\eu^{\,n\ln(1 + 1/n)} = \eu^{\,1 - \frac1{2n} + O(n^{-2})} =
\eu\bigl(1 - \frac{1}{2n} + O(n^{-2})\bigr)$، ومنه
\[
\eu - \Bigl(1 + \frac1n\Bigr)^{\!n} \sim \frac{\eu}{2n} :
\]
وهي حدود موجبة مكافئة لمضاعف توافقي: فتتباعد.
\end{solution}

\begin{solution}{exo:b2:series:12}
بالرتابة، $n\,a_{2n} \leq a_{n+1} + a_{n+2} + \dots +
a_{2n} = S_{2n} - S_n \to 0$ (بمحك كوشي للمتسلسلة
المتقاربة). ومنه $2n\,a_{2n} \to 0$، و$(2n{+}1)\,
a_{2n+1} \leq (2n{+}1)a_{2n} = \frac{2n+1}{2n}\,(2n\,a_{2n}) \to
0$: فتؤول المتتاليتان الجزئيتان من $(na_n)$ إلى $0$، ومنه $na_n \to 0$.

\emph{ويفشل العكس:} فالمتتالية $a_n = \frac1{n\ln n}$ موجبة
متناقصة مع $na_n = \frac1{\ln n} \to 0$، ومع ذلك تتباعد $\sum a_n$
(حدود برتران، \cref{pb:b2:series:1}، السؤال
18). \emph{والرتابة جوهرية:} فلتكن $a_n = \frac1n$ حين يكون $n$
قوة للعدد $2$، و$a_n = 2^{-n}$ فيما عدا ذلك: عندئذٍ $\sum a_n \leq
\sum_k 2^{-k} + \sum_n 2^{-n} < \infty$، لكن $na_n = 1$ على
قوى العدد $2$: أي $na_n \not\to 0$.
\end{solution}

\begin{solution}{pb:b2:series:1}
\textbf{1.} بالتراجع على $m$: من أجل $m = 0$ يساوي الطرفان $1$؛
وتضرب الخطوة في $\cos\theta + \iu\sin\theta$ وتستعمل
صيغتَي الجمع $\cos(m\theta + \theta) =
\cos m\theta\cos\theta - \sin m\theta\sin\theta$ و$\sin(m\theta
+ \theta) = \sin m\theta\cos\theta + \cos m\theta\sin\theta$.
وبالنشر بدل ذلك بمبرهنة ثنائي الحدّ مع $m = 2n+1$
وجمع الجزء التخيلي (أي القوى الفردية للمقدار
$\iu\sin\theta$، مع $\iu^{2j+1} = (-1)^j\iu$):
\[
\sin\bigl((2n{+}1)\theta\bigr) = \sum_{j=0}^{n}(-1)^j
\binom{2n+1}{2j+1}\cos^{2(n-j)}\theta\,\sin^{2j+1}\theta .
\]

\textbf{2.} على $\intoo0{\frac\pi2}$، $\sin\theta \neq 0$: فنُخرِج
$\sin^{2n+1}\theta$ من كل حدّ، فيبقى
$\bigl(\frac{\cos^2\theta}{\sin^2\theta}\bigr)^{n-j} =
(\cot^2\theta)^{n-j}$: أي المتطابقة المعروضة مع $P_n(x) =
\sum_j(-1)^j\binom{2n+1}{2j+1}x^{n-j}$. ومعاملها من الدرجة $n$
هو معامل $j = 0$ أي $\binom{2n+1}{1} = 2n + 1 \neq 0$.

\textbf{3.} عند $\theta_k = \frac{k\pi}{2n+1}$:
$\sin\bigl((2n{+}1)\theta_k\bigr) = \sin k\pi = 0$ في حين أن
$\sin^{2n+1}\theta_k \neq 0$، ومنه $P_n(\cot^2\theta_k) = 0$. وتتزايد
$\theta_k$ تمامًا في $\intoo0{\frac\pi2}$، حيث تتناقص
$\cot^2$ تمامًا: فالقيم $x_k =
\cot^2\theta_k$ مختلفة مثنى مثنى --- أي $n$ جذرًا مختلفًا
لكثير حدود درجته $n$، ومن ثم فهي كل جذوره.

\textbf{4.} فييت: مجموع الجذور هو سالب نسبة
معاملَي $x^{n-1}$ و$x^n$:
\[
\sum_{k=1}^{n}\cot^2\theta_k =
\frac{\binom{2n+1}{3}}{\binom{2n+1}{1}}
= \frac{(2n+1)(2n)(2n-1)/6}{2n+1} = \frac{n(2n-1)}{3}.
\]

\textbf{5.} $\frac{1}{\sin^2\theta} = 1 + \cot^2\theta$:
وبالجمع، $n + \frac{n(2n-1)}3 = \frac{3n + 2n^2 - n}{3} =
\frac{2n(n+1)}{3}$.

\textbf{6.} على $\intoo{0}{\frac\pi2}$: $\sin\theta < \theta <
\tan\theta$ (مجلد السنة الأولى). وأخذ المقلوبات يعكس:
$\cot\theta < \frac1\theta < \frac1{\sin\theta}$، والتربيع
(وكلها موجبة) يعطي $\cot^2\theta < \frac1{\theta^2} <
\frac1{\sin^2\theta}$.

\textbf{7.} نجمع السؤال 6 عند $\theta = \theta_k$ على $k \leq
n$، مستعملين السؤالين 4 و5، و$\frac1{\theta_k^2} =
\frac{(2n+1)^2}{k^2\pi^2}$:
\[
\frac{n(2n-1)}3 < \frac{(2n+1)^2}{\pi^2}\sum_{k=1}^n\frac1{k^2}
< \frac{2n(n+1)}3 .
\]

\textbf{8.} نضرب في $\frac{\pi^2}{(2n+1)^2}$:
\[
\frac{\pi^2}{3}\cdot\frac{n(2n-1)}{(2n+1)^2}
< \sum_{k=1}^{n}\frac1{k^2} <
\frac{\pi^2}{3}\cdot\frac{2n(n+1)}{(2n+1)^2}.
\]
ويؤول الحدّان كلاهما إلى $\frac{\pi^2}3\cdot\frac12 = \frac{\pi^2}6$
(فالكسور الناطقة تؤول إلى $\frac12$). وتتزايد المجاميع الجزئية،
ومن ثم فهي متقاربة، ويعطي الحصر $\zeta(2) =
\frac{\pi^2}6$: وهي مبرهنة أويلر، ببرهان كوشي.

\textbf{9.} تتزايد المجاميع الجزئية إلى $\zeta(2) =
\frac{\pi^2}6$، ومنه $0 \leq \frac{\pi^2}6 - \sum_{k\leq n}k^{-2}$؛
ويعطي الحدّ الأدنى في السؤال 8 أن
\[
\frac{\pi^2}6 - \sum_{k\leq n}\frac1{k^2}
\leq \frac{\pi^2}6 - \frac{\pi^2}3\cdot\frac{n(2n-1)}{(2n+1)^2}
= \frac{\pi^2}6\cdot\frac{(2n+1)^2 - (4n^2 -
2n)}{(2n+1)^2}
= \frac{\pi^2}6\cdot\frac{6n + 1}{(2n+1)^2}
= O\Bigl(\frac1n\Bigr),
\]
وهو يطابق الذيل الدقيق $\frac1n - \frac1{2n^2} + O(n^{-3})$ في
\cref{exo:b2:comparison:11}.

\textbf{10.} دالة فييت المتناظرة الأولية الثانية
للجذور هي $\sigma_2 = \frac{\binom{2n+1}5}{\binom{2n+1}1} =
\frac{(2n)(2n-1)(2n-2)(2n-3)}{120}$، ومن ثم
\[
\sum_k\cot^4\theta_k = \sigma_1^2 - 2\sigma_2
= \Bigl(\frac{n(2n-1)}3\Bigr)^2 -
\frac{2n(2n-1)(2n-2)(2n-3)}{60}
\sim \frac{4n^4}9 - \frac{4n^4}{15} = \frac{8n^4}{45}.
\]
وبحصر $\cot^4\theta < \theta^{-4} < (1 + \cot^2\theta)^2 =
1 + 2\cot^2\theta + \cot^4\theta$ والجمع: يكون المجموعان الخارجيان
$\frac{8n^4}{45}(1 + o(1))$ (فالمقدار المضاف $n + 2\sigma_1 =
O(n^2)$ مهمَل)، في حين أن الأوسط
$\frac{(2n+1)^4}{\pi^4}\sum_{k\leq n}k^{-4}$. ومنه
\[
\sum_{k\leq n}\frac1{k^4} \longrightarrow
\pi^4\cdot\frac{8/45}{16} = \frac{\pi^4}{90}.
\]

\textbf{11.} بفصل $\zeta(2)$ حسب التماثل: $\sum_{\text{زوجي}}
= \sum_j\frac{1}{(2j)^2} = \frac14\zeta(2) = \frac{\pi^2}{24}$،
ومنه $\sum_{\text{فردي}} = \zeta(2) - \frac{\pi^2}{24} =
\frac{\pi^2}8$. والمتناوبة: $\sum_k\frac{(-1)^{k-1}}{k^2} =
\sum_{\text{فردي}} - \sum_{\text{زوجي}} = \frac{\pi^2}8 -
\frac{\pi^2}{24} = \frac{\pi^2}{12}$ (ويبرّر التقارب المطلق
إعادة التجميع،
\cref{thm:b2:series:rearrangement}).

\textbf{12.} $S = \frac{\zeta(2)^2}{\zeta(4)} =
\frac{(\pi^2/6)^2}{\pi^4/90} = \frac{90}{36} = \frac52$.
والحدس: تقول المتطابقة $\zeta(2)^2 = \zeta(4)S$ في
\cref{exo:b2:series:7} إن إخراج العامل $\gcd$
يُنظّم الأزواج إلى أزواج أولية فيما بينها؛ ومقلوب المقدار
$\frac1{\zeta(2)} = \frac6{\pi^2} \approx 0.608$ هو المرشح الطبيعي
لكثافة الأزواج الأولية فيما بينها بين كل الأزواج ---
وهو قول عن $\lim_N \frac{1}{N^2}\#\{(m,n) \leq N :
\gcd = 1\}$ برهانه الأمين (بحدود الخطأ) ينتمي إلى
مجلد السنة الثالثة.

\textbf{13.} خطأ الجمع المباشر $\sim \frac1n$: فستة
أرقام تقتضي نحو $10^6$ حدّ. أما المجموع المصحَّح
$\sum_{k\leq n}k^{-2} + \frac1n - \frac1{2n^2}$ فخطؤه
$O(n^{-3})$: وعند $n = 100$ يساوي $1.6449339\dots$ في مقابل
$\frac{\pi^2}6 = 1.6449341\dots$ --- أي خطأ $1.7\cdot10^{-7}$،
فسبعة أرقام من مئة حدّ. فالتصحيحات المقاربة تغلب
الصبر الخام بأربع رتب من الحجم.

\textbf{14.} $n = 1$: $P_1(x) = 3x - 1$، والجذر $\frac13$، وفعلًا
$\cot^2\frac\pi3 = \bigl(\frac1{\sqrt3}\bigr)^2 =
\frac13 = \frac{1\cdot1}3$. و$n = 2$: تتنبأ الصيغة بالقيمة
$\frac{2\cdot3}3 = 2$؛ ومع $\cos\frac\pi5 = \frac{1 +
\sqrt5}{4}$، نحسب $\cot^2 36^\circ \approx 1.894$ و
$\cot^2 72^\circ \approx 0.106$: فالمجموع $2.000$.

\textbf{15.} الأعداد $\tan^2\theta_k = \frac1{x_k}$ هي
جذور $Q(x) = x^nP_n\bigl(\frac1x\bigr) =
\sum_{j=0}^n(-1)^j\binom{2n+1}{2j+1}x^j$ (لأن الأعداد $x_k$
غير معدومة). وفييت على $Q$: المعامل الرئيسي هو $(-1)^n$
(الحدّ $j = n$)، والتالي هو $(-1)^{n-1}\binom{2n+1}{2n-1} =
(-1)^{n-1}\binom{2n+1}{2}$، ومنه
\[
\sum_{k=1}^n\tan^2\theta_k =
-\frac{(-1)^{n-1}\binom{2n+1}2}{(-1)^n} = \binom{2n+1}2\cdot
\frac{2}{2n+1}\cdot\frac{2n+1}{2}
= n(2n+1).
\]
وللتحقق عند $n = 1$: $\tan^2\frac\pi3 = 3 = 1\cdot3$.

\textbf{16.} ليكن $\gamma = \frac{1+\alpha}2 \in
\intoo{1}{\alpha}$. عندئذٍ $\frac{n^{-\alpha}(\ln
n)^{-\beta}}{n^{-\gamma}} = n^{\gamma - \alpha}(\ln n)^{-\beta}
\to 0$ (فقوة سالبة للمقدار $n$ تغلب أي قوة للمقدار $\ln n$)، ومن ثم
تصير الحدود ابتداءً من رتبة ما $\leq n^{-\gamma}$ مع $\gamma > 1$:
فالتقارب بالمقارنة مع متسلسلة ريمان.

\textbf{17.} ليكن $\gamma = \frac{1+\alpha}2 \in
\intoo{\alpha}{1}$: فالآن $\frac{n^{-\gamma}}{n^{-\alpha}(\ln
n)^{-\beta}} = n^{\alpha-\gamma}(\ln n)^{\beta} \to 0$، ومن ثم
تصير الحدود ابتداءً من رتبة ما $\geq n^{-\gamma}$ مع $\gamma < 1$:
فالتباعد.

\textbf{18.} $f(t) = \frac1{t(\ln t)^\beta}$ موجبة
ومتصلة ومتناقصة من أجل $t$ الكبيرة (فمشتق لوغاريتمها
$-\frac1t\bigl(1 + \frac{\beta}{\ln t}\bigr) < 0$
ابتداءً من رتبة ما). والدوال الأصلية: من أجل $\beta \neq 1$،
$\int^x f = \frac{(\ln x)^{1-\beta}}{1-\beta} + \text{ثابت}$،
ولها نهاية منتهية إذا وفقط إذا $\beta > 1$؛ ومن أجل $\beta = 1$،
$\int^x f = \ln\ln x \to \infty$. وحسب
\cref{thm:b2:comparison:seriesintegral}، للمتسلسلة والتكامل
الطبيعة نفسها: فالتقارب إذا وفقط إذا $\beta > 1$.

\textbf{19.} المحك نفسه: $\frac{\dd}{\dd t}\ln\ln\ln t =
\frac{1}{t\ln t\,\ln\ln t}$، و$\ln\ln\ln t \to \infty$:
فالتباعد. و$\frac{\dd}{\dd t}\Bigl(-\frac1{\ln\ln t}\Bigr)
= \frac{1}{t\ln t\,(\ln\ln t)^2}$ مع $-\frac1{\ln\ln t} \to
0$: فالتقارب.

\textbf{20.} الأولى: $n^{1/\ln n} = \eu^{\ln n/\ln n} = \eu$، ومن ثم
تكون الحدود بالضبط $\frac{1}{\eu\,n}$: أي مضاعف للمتسلسلة
التوافقية، فهي \emph{متباعدة} --- إذ إن الأس $1 +
\frac1{\ln n}$ يزحف إلى $1$ بسرعة زائدة. والثانية: $n^{1/\ln\ln n}
= \eu^{\ln n/\ln\ln n}$، و$\frac{\ln n}{\ln\ln n} \geq
2\ln\ln n$ ابتداءً من رتبة ما، ومنه $n^{1/\ln\ln n} \geq (\ln n)^2$:
فالحدود $\leq \frac1{n(\ln n)^2}$، وهي متسلسلة برتران متقاربة
(السؤال 18): فهي \emph{متقاربة}. وتمرّ الحدود
تمامًا بين هذين الأسّين.

\textbf{21.} $R_n \downarrow 0$ و
\[
\sqrt{R_n} - \sqrt{R_{n+1}} = \frac{R_n -
R_{n+1}}{\sqrt{R_n} + \sqrt{R_{n+1}}} =
\frac{a_n}{\sqrt{R_n} + \sqrt{R_{n+1}}} \geq
\frac{a_n}{2\sqrt{R_n}},
\]
ومنه $\sum_n \frac{a_n}{\sqrt{R_n}} \leq 2\sum_n(\sqrt{R_n} -
\sqrt{R_{n+1}}) = 2\sqrt{R_1} < \infty$ (بالتلسكوب). ومع ذلك
$\frac{a_n/\sqrt{R_n}}{a_n} = \frac1{\sqrt{R_n}} \to \infty$:
فالمتسلسلة الجديدة تتقارب وهي أكبر بما لا نهاية. فلا توجد
متسلسلة متقاربة أبطأ من الجميع؛ ولا يمكن لمحكّات المقارنة
مع أي عائلة ثابتة أن تكون تامة أبدًا.

\textbf{22.} من أجل $(a_n)$ موجبة متناقصة، نجمّع الحدود
بين قوتين متتاليتين للعدد $2$:
\[
2^{k}a_{2^{k+1}} \leq \sum_{n=2^k}^{2^{k+1}-1} a_n \leq
2^ka_{2^k} .
\]
وبالجمع على $k$: إذا تقاربت $\sum 2^ka_{2^k}$، تكون المجاميع
الجزئية للمتسلسلة $\sum a_n$ محدودة (فتتقارب)؛ وإذا تقاربت $\sum a_n$،
فإن $\sum_k 2^{k+1}a_{2^{k+1}} \leq 2\sum_n a_n <
\infty$. ومن أجل $a_n = \frac1{n(\ln n)^\beta}$: $2^ka_{2^k} =
\frac{1}{(k\ln 2)^\beta}$، وتتقارب $\sum k^{-\beta}$ إذا وفقط إذا
$\beta > 1$: أي حدود السؤال 18 من جديد، بلا
تكاملات.

\textbf{23.} بمقارنة التكامل،
\[
\sum_{n > N}\frac{1}{n(\ln n)^2} \geq
\int_{N+1}^{\infty}\frac{\dd t}{t(\ln t)^2} =
\frac{1}{\ln(N+1)},
\]
وهو عند $N = 10^6$ يساوي $\approx 0.0724$: فبعد مليون حدّ
يبقى الذيل أكبر من $0.07$ --- فالمتسلسلة تتقارب، لكن لا
جمع مباشر يُظهر مجموعها أبدًا. وقارِن ذلك بالسؤال
13، حيث اشترى تصحيح مقارب واحد سبعة
أرقام من مئة حدّ: فمعرفة \emph{كيف} تتقارب متسلسلة
أثمن من معرفة أنها تتقارب.

\textbf{24.} $\sum\frac1{n\ln n}$: تتباعد ($\alpha = 1$،
$\beta = 1$، السؤال 18). و$\sum\frac1{n^{1.01}}$: تتقارب
(ريمان، $\alpha > 1$). و$\sum\frac{(\ln n)^{100}}{n^{1.001}}$:
تتقارب ($\alpha = 1.001 > 1$، $\beta = -100$، السؤال 16).
و$\sum\frac1{n\ln n(\ln\ln n)^3}$: تتقارب (بنمط السؤال 19:
فالدالة الأصلية $-\frac12(\ln\ln t)^{-2}$، ولها نهاية
منتهية).

\textbf{25.} يحوّل دي موافر انعدام $\sin(2n{+}1)
\theta_k$ إلى انعدام كثير حدود عند
$\cot^2\theta_k$، ويقرأ فييت مجاميع الجذور الدقيقة التي لم يكن
التحليل وحده ليقدّمها إلا تقديرًا (الأسئلة 1--5). واحتاج
الحصر إلى القيمتين \emph{الدقيقتين} $\frac{n(2n-1)}3$ و
$\frac{2n(n+1)}3$ على الطرفين --- فالمكافئات كانت
ستصادر المطلوب، إذ إن كل المغزى هو الثابت
$\frac{\pi^2}6$ (السؤالان 7--8). وجرى الجزء الرابع كله على
مقارنة المتسلسلة بالتكامل في \cref{ch:b2:comparison}، وقامت
الدوال الأصلية اللوغاريتمية بالتصنيف (السؤالان
18--19). ويهدم السؤال 21 حلمَ محك المقارنة
الشامل: فتحت كل متسلسلة متقاربة تقع أخرى
أبطأ منها بما لا نهاية --- والسلالم مثل خريطة برتران ترسم الحدود
بدقة متزايدة أبدًا لكنها لا تبلغها. والقمتان: مقدار أويلر $\zeta(2)
= \frac{\pi^2}6$ مع طابقه الأعلى $\zeta(4) =
\frac{\pi^4}{90}$ (السؤالان 8 و10)، وتصنيف
برتران (الأسئلة 16--18)؛ ويعود $\zeta(2)$ في
\cref{ch:b2:fourier}، حيث تعيد متطابقة بارسفال برهانه
في سطر واحد انطلاقًا من متسلسلة فورييه للموجة المسننة --- ثابت
واحد، وحضارتان.
\end{solution}
